\documentclass[11pt,a4paper]{article}

\usepackage[utf8]{inputenc}
\usepackage[T1]{fontenc}
\usepackage{lmodern}
\usepackage[margin=1in]{geometry}
\usepackage{amsmath,amssymb,amsthm,mathtools}
\usepackage{algorithm}
\usepackage{algpseudocode}
\usepackage{booktabs}
\usepackage[numbers,sort&compress]{natbib}
\usepackage[hidelinks]{hyperref}
\usepackage{cleveref}
\usepackage{microtype}
\setlength{\emergencystretch}{3em}

\theoremstyle{plain}
\newtheorem{theorem}{Theorem}[section]
\newtheorem{proposition}[theorem]{Proposition}
\newtheorem{corollary}[theorem]{Corollary}
\theoremstyle{definition}
\newtheorem{definition}[theorem]{Definition}
\theoremstyle{remark}
\newtheorem{remark}[theorem]{Remark}

\newcommand{\R}{\mathbb{R}}
\newcommand{\N}{\mathbb{N}}
\newcommand{\norm}[1]{\left\lVert #1 \right\rVert}
\newcommand{\abs}[1]{\left\lvert #1 \right\rvert}
\newcommand{\ip}[2]{\left\langle #1, #2 \right\rangle}
\newcommand{\lean}[1]{\path{#1}}
\DeclareMathOperator{\Var}{Var}
\DeclareMathOperator{\Cov}{Cov}
\DeclareMathOperator{\loo}{loo}
\DeclareMathOperator{\Sh}{Sh}
\DeclareMathOperator{\KL}{KL}

\title{Scale and Attribution in an Evolving Knowledge Store:\\
When Macro Behaviour Can Be Studied,\\
and How Progress Is Attributed to Sources}
\author{Alberto Serrano-Calva}
\date{November 7, 2026}

\begin{document}
\maketitle

\begin{abstract}
A store that builds context from many sources should be able to say, with evidence, which
sources, context windows and analysis tools improve its results. That needs data at a scale where
the macro quantities of the store's graph are stable, and a continuous evaluation whose outcomes
can be attributed to each source or tool. We answer both with theorems machine-checked in Lean~4
over Mathlib before any code relies on them. Under the lazy walk $1-\varepsilon L$ on a connected
graph with spectral gap $\gamma$, the time average over $T$ steps of a mean-zero observable has
energy at most $\norm{x}^2/(T\varepsilon\gamma)^2$, so a window must span many relaxation times $1/(\varepsilon\gamma)$, and
the sampled form gives at least $T\varepsilon\gamma/2$ effective samples. With an eigengap $g$ between the
kept and the dropped modes, the dropped part's share of the energy falls as $e^{-2gt}$ under the
heat kernel, so a $k$-mode description explains all but $\eta$ of what remains once
$t \ge \log(R/\eta)/(2g)$; a refuted control shows that without a gap no time does. The Shapley
value, the average over orderings of a source's marginal contribution, is efficient, symmetric,
null-respecting and additive; it is the only attribution with the four, by the decomposition of
every score into unanimity games with Harsanyi dividends; and leave-one-out, which an ablation arm
measures, agrees with it on constant marginals and confirms neither of two redundant sources. The
number of resolved units an effect needs is the ledger's power calculation, derived from the normal
approximation. We design the loop in which the store's own future labels its present, forward only
and pre-registered, with each configuration an arm and the log-score advantage over a reference arm
as the common currency, and we list the open corpora that cross the thresholds. Each theorem is
attacked on an exact rational shadow with a refuted control, and the library is replayed through
Lean's kernel and five independent checkers.
\end{abstract}

\section{Introduction}

A knowledge store that is kept over time is a sequence of graphs: notes are written and rewritten,
citations are added and severed, sources arrive. The previous article in this series
(\lean{papers/2026-11-07-graph-dynamics}) proved the dynamics of such a graph: the sum over walks
is a resolvent, a flow is read
through its Hodge decomposition, and the macro regime is the span of the low modes of the
Laplacian, which evolve exactly by their own eigenvalues while the rest decays at the rate of the
spectral gap. This article asks two questions that follow from it.

The first is \emph{scale}. The macro regime is a theorem on a fixed graph, but the quantities it is
read from, the spectral gap $\gamma$, the eigengaps between modes, a time average of an observable,
are measured on a sample of a growing corpus, and on a small sample they fluctuate with every note
added. At what scale do they stop fluctuating, and how long a window does an estimate of a macro
observable need? We prove two thresholds (the relaxation time for a time average, the eigengap at
the cut for a macro description), cite a third (the concentration of random graph spectra once the
minimum degree exceeds a constant times $\log N$), and say plainly how far below every threshold a
small private corpus is.

The second is \emph{attribution}. A store should be able to say which of its sources, which
context windows and which analysis tools improve its results, and say it with evidence rather than
by assertion. We design a continuous evaluation loop: at each cut of the store's history, questions
about events after the cut (a link that will form, a note that will be superseded, a version that
will appear, an answer a later source will confirm) are registered in a pre-registered prediction
ledger, forecast by every configuration of the system, and resolved by a later cut, forward only.
Each configuration is an arm, and the outcome is attributed by two readings: a randomization test
per ablation, which measures leave-one-out, and the Shapley value over the sources. We prove the
Shapley value's existence and uniqueness from its four axioms, prove when leave-one-out agrees with
it, and exhibit the case where it does not: two redundant sources, either sufficient, each
attributed nothing by leave-one-out and half by Shapley. The disagreement is not a defect of either
reading; it is the finding that the store holds the same evidence twice.

The contributions are:
\begin{enumerate}
\item SCL-1, the estimation threshold (\Cref{thm:scl1}): under the lazy walk on a connected graph
  with gap $\gamma$, the time average of a mean-zero field over $T$ steps has energy at most
  $((1/T)\sum_{n<T}(1-\varepsilon\gamma)^n)^2\norm{x}^2 \le \norm{x}^2/(T\varepsilon\gamma)^2$, and
  precision $\delta$ needs $T \ge \norm{x}/(\varepsilon\gamma\delta)$. The sampled form, derived,
  bounds the variance of the time average by $2\sigma^2/(\varepsilon\gamma T)$, at least $T\varepsilon\gamma/2$ effective samples.
\item SCL-2, scale separation (\Cref{thm:scl2}): the dropped modes' energy relative to the kept
  modes' falls by $e^{-2(\mu-\nu)t}$, with the time $\log(R/\eta)/(2g)$ at which a $k$-mode
  description explains all but $\eta$, a discrete twin under the lazy walk, and the refuted
  control without a gap.
\item SCL-3, attribution (\Cref{thm:scl3}): the Shapley value satisfies the four axioms, is the
  unique attribution that does, equals the Harsanyi dividend formula, and leave-one-out agrees with
  it exactly on scores with constant marginals, with a witness for the general disagreement.
\item SCL-4, power (\Cref{sec:scl4}): the number of independent resolved units to detect an effect,
  derived from the normal approximation in the form the ledger uses.
\item The loop (\Cref{sec:loop}), the datasets (\Cref{sec:datasets}) and the recipe
  (\Cref{sec:recipe}), each tied to the thresholds and the theorems.
\end{enumerate}
Each theorem is stated with its Lean name in \lean{formal/Palace/Scale.lean}, attacked on an exact
rational shadow with a negative control in \lean{formal/PalaceTest/Scale.lean}, and replayed, with
the rest of the library, through Lean's kernel and five independent checkers.

\section{Background and related work}

\paragraph{Spectral gaps and mixing.}
For a weighted graph with Laplacian $L = D - A$ on $N$ nodes, the spectral gap is the least
nonzero eigenvalue on a connected support, and the lazy walk $1-\varepsilon L$ with
$\varepsilon \le 1/(2d_{\max})$, so that $\varepsilon\lambda \le 1$ on the whole spectrum since $\lambda_{\max} \le 2d_{\max}$, is a reversible Markov chain whose relaxation time is $1/(\varepsilon\gamma)$
\cite{levin2017mixing,chung1997spectral,lovasz1993random}. Its mixing is geometric at the gap's
rate, and the variance of a time average along the chain is governed by the same gap: for a chain
started at stationarity the asymptotic variance of an observable is at most $(2/\gamma)$ times its
stationary variance \cite{aldous2002reversible,levin2017mixing}. SCL-1 proves the deterministic
form of this bound for the exact diffusion and derives the sampled form from it.

\paragraph{Concentration of random graph spectra.}
For a random graph with independent edges whose minimum expected degree $d_{\min}$ is at least a constant
times $\log N$, the adjacency matrix and the normalized Laplacian concentrate around their
expectations in operator norm, within $O(\sqrt{\log N/d_{\min}})$ with high probability
\cite{oliveira2009concentration}, with explicit constants through a matrix Chernoff bound
\cite{chung2011spectra}. This is the threshold at which a measured spectrum is the corpus's rather
than the sample's; it is cited here, not proved, and a sibling study measures its shape on real
data.

\paragraph{Coarse-graining and scale separation.}
A function of a symmetric operator acts on each eigenvector by the function of its eigenvalue, so
the heat kernel commutes with the projection onto any set of modes and the dropped part decays at
the rate of the gap below it; that is DYN-4 of the previous article. Scale separation in the sense
used here, that a few slow modes describe the system on a time scale set by the gap to the fast
ones, is the standard adiabatic picture; SCL-2 states it as an inequality between the two parts'
energies with both bounds proved.

\paragraph{Data valuation.}
The Shapley value \cite{shapley1953value} is the unique attribution satisfying efficiency,
symmetry, the null player axiom and additivity, and its decomposition through the unanimity games
and their dividends is Harsanyi's \cite{harsanyi1963bargaining}. Its use to value data for a learned
model, with the value function the model's score on a subset of the data, is Data Shapley
\cite{ghorbani2019data}, estimated by Monte Carlo over permutations \cite{castro2009polynomial},
with more efficient estimators under structural assumptions \cite{jia2019efficient}. SCL-3 proves
the axioms and the uniqueness in Lean for a general value function and adds the comparison with
leave-one-out that an ablation measures.

\paragraph{Ablation and pre-registration.}
An ablation removes one component and measures the drop; it measures necessity, and SCL-3 says
exactly when that equals contribution. The evaluation here runs under a protocol written before
the first prediction \cite{nosek2018preregistration}, scored by strictly proper rules
\cite{brier1950verification,gneiting2007strictly}, with the family-wise error controlled across
secondary comparisons \cite{holm1979simple} and the inference at the cluster level
\cite{cameron2008bootstrap}; the forward-only rule follows the practice of keeping only questions
whose outcome a model could not have seen \cite{halawi2024approaching}. The log-score advantage over
a reference arm is the growth rate of a proportional stake \cite{kelly1956}, the identity the
ledger's growth readout rests on.

\paragraph{Temporal self-supervision and link prediction.}
That the store's own future can label its present is the link-prediction setting
\cite{libennowell2007link}: a snapshot as of a cut, and the question of which links appear by a
later cut, with the snapshot's own structure as the only input. The loop generalizes it to
supersessions, new versions and confirmations by later sources, each resolved by the store at a
later cut.

\section{Theory}

Throughout, $V$ is a finite set of $N$ nodes, $A$ a symmetric nonnegative weight matrix on $V$
with Laplacian $L = D - A$, $M$ a real symmetric matrix with orthonormal eigenvectors $\varphi_j$
and eigenvalues $\lambda_j$ (the \emph{modes}, \lean{Palace.Dynamics.mode}), $H_t = e^{-tM}$ its
heat kernel by the functional calculus, and $P_S f = \sum_{j \in S}\ip{f}{\varphi_j}\varphi_j$ the
projection of a field onto a set $S$ of modes. The spectral gap is stated as in DYN-6: $\gamma$ is a
gap of $A$ when $\gamma \le \lambda_j$ for every nonzero eigenvalue of $L$, which on a connected
support is the Poincar\'e inequality $\gamma\norm{f}^2 \le f^{\top}Lf$ for every $f \perp \mathbf{1}$.

\subsection{SCL-1: the estimation threshold}

\begin{definition}[time average]
For an operator $N$ on $V$, a window of $T$ steps and a field $x$, the time average is
$\operatorname{avg}_T(x) = \frac{1}{T}\sum_{n<T} N^n x$ (\lean{Palace.Scale.timeAverage}): the
estimate of a macro observable from $T$ consecutive readings of the diffusion started at $x$.
\end{definition}

\begin{theorem}[SCL-1, the estimation threshold]\label{thm:scl1}
Let $A$ have a connected support, let $\gamma$ be a gap of $A$, and let $\varepsilon \ge 0$ with
$\varepsilon\lambda_j \le 1$ on every eigenvalue of $L$. For every $x \perp \mathbf{1}$ and every
window $T$, under the lazy walk $N_\varepsilon = 1 - \varepsilon L$,
\[
  \norm{\operatorname{avg}_T(x)}^2
  \;\le\; \Bigl(\frac{1}{T}\sum_{n<T}(1-\varepsilon\gamma)^n\Bigr)^2 \norm{x}^2
  \quad\text{(\lean{lazy\_timeAverage\_energy\_le})},
\]
and when $0 < \varepsilon\gamma \le 1$ and $T \ge 1$,
\[
  \norm{\operatorname{avg}_T(x)}^2 \;\le\; \frac{\norm{x}^2}{(T\varepsilon\gamma)^2}
  \quad\text{(\lean{lazy\_timeAverage\_energy\_le\_window})}.
\]
Consequently, for a precision $\delta > 0$, $T \ge \norm{x}/(\varepsilon\gamma\delta)$ implies
$\norm{\operatorname{avg}_T(x)}^2 \le \delta^2$ (\lean{lazy\_timeAverage\_precision}).
\end{theorem}

\begin{proof}[Proof sketch]
On a mode $\varphi_j$ the lazy walk multiplies by $g_j = 1 - \varepsilon\lambda_j$, so the time
average multiplies by the geometric mean sum $\frac{1}{T}\sum_{n<T} g_j^n$
(\lean{timeAverage\_dotProduct\_mode}, through the symmetry of $N^n$). By Parseval,
$\norm{\operatorname{avg}_T(x)}^2 = \sum_j (\frac{1}{T}\sum_{n<T} g_j^n)^2\ip{x}{\varphi_j}^2$. A
mean-zero field on a connected support has no weight on the zero mode (DYN-5's
\lean{dotProduct\_mode\_eq\_zero}), so on every mode it lives on $0 \le g_j \le 1 - \varepsilon\gamma$,
and the geometric sums are monotone in their ratio (\lean{timeAverage\_energy\_le}). The second
bound is the geometric sum under $1/(1-\rho) = 1/(\varepsilon\gamma)$ (\lean{sum\_pow\_le\_inv\_one\_sub},
Mathlib's \lean{geom\_sum\_Ico\_le\_of\_lt\_one}), and the precision is its square root.
\end{proof}

In continuous units a window of $T$ lazy steps spans $T_c = T\varepsilon$ of diffusion time and
$1/\gamma$ is the relaxation time, so the bound reads $\norm{x}^2/(\gamma T_c)^2$: the estimate is
within $\delta$ of the mean once the window spans $\norm{x}/\delta$ relaxation times. This is the
deterministic form, in which the readings are the exact states of the diffusion, and its constant
$1/(\gamma T_c)^2$ is the one proved.

\begin{proposition}[the sampled form, derived]\label{prop:sampled}
Let $X_0, X_1, \dots$ be the lazy walk started from its stationary distribution (uniform, since
$1-\varepsilon L$ is symmetric and stochastic for $\varepsilon \le 1/d_{\max}$; take $\varepsilon \le 1/(2d_{\max})$, which also gives $\varepsilon\lambda_j \le 1$ as in \Cref{thm:scl1}), and $f$ an observable with
centered part $f_0$ and stationary variance $\sigma^2 = \norm{f_0}^2/N$. Then
\[
  \Var\Bigl(\frac{1}{T}\sum_{t<T} f(X_t)\Bigr)
  \;\le\; \frac{\sigma^2}{T}\cdot\frac{2-\varepsilon\gamma}{\varepsilon\gamma}
  \;\le\; \frac{2\sigma^2}{\varepsilon\gamma T},
\]
so $T \ge 2\sigma^2/(\varepsilon\gamma\delta^2)$ lazy steps, that is $T_c = \varepsilon T \ge 2\sigma^2/(\gamma\delta^2)$ of diffusion time, suffice for a standard error of at most $\delta$, and
a window carries at least $T\varepsilon\gamma/2 = T_c\gamma/2$ effective samples, $\sigma^2$ over the variance of the average.
\end{proposition}

\begin{proof}[Derivation]
$\Var(\sum_t f(X_t)) = \sum_{s,t}\Cov(f(X_s), f(X_t)) = \sum_{s,t}\ip{f_0}{P^{\abs{s-t}} f_0}/N$ for
a stationary reversible chain with transition matrix $P = 1 - \varepsilon L$, and each term is at most
$(1-\varepsilon\gamma)^{\abs{s-t}}\norm{f_0}^2/N$ by Cauchy--Schwarz and DYN-6's contraction on the
modes (\lean{lazy\_pow\_energy\_le}). Summing the geometric series in $\abs{s-t}$ gives
$T\sigma^2(1 + 2\sum_{k \ge 1}(1-\varepsilon\gamma)^k) = T\sigma^2(2-\varepsilon\gamma)/(\varepsilon\gamma)$. The
contraction is machine-checked; the chain, its stationarity and the covariance identity are not
formalized, which is the textbook part \cite{aldous2002reversible,levin2017mixing}.
\end{proof}

Both forms say the same thing: a macro observable cannot be estimated from a window shorter than
the relaxation time, whatever the number of readings in it, and the number of readings that count
is the window measured in relaxation times.

\subsection{SCL-2: scale separation}

DYN-4 bounds the dropped part from above: for $\mu$ at most every dropped eigenvalue and $t \ge 0$,
$\norm{H_t(f - P_Sf)}^2 \le e^{-2\mu t}\norm{f - P_Sf}^2$ (\lean{heat\_sub\_modeProj\_le}). Scale
separation needs the kept part bounded from below.

\begin{theorem}[SCL-2, scale separation]\label{thm:scl2}
Let $S$ be a set of modes of $M$, $\mu \le \lambda_j$ for every $j \notin S$ and $\lambda_j \le \nu$
for every $j \in S$. Then for $t \ge 0$,
\[
  e^{-2\nu t}\norm{P_Sf}^2 \;\le\; \norm{H_t P_S f}^2
  \quad\text{(\lean{heat\_modeProj\_energy\_ge})},
\]
and
\[
  \norm{H_t(f - P_Sf)}^2\,\norm{P_Sf}^2
  \;\le\; e^{-2(\mu-\nu)t}\,\norm{f - P_Sf}^2\,\norm{H_tP_Sf}^2
  \quad\text{(\lean{scale\_separation})}.
\]
If moreover $g = \mu - \nu > 0$, $\eta > 0$, $\norm{P_Sf} > 0$ and
$t \ge \log(R/\eta)/(2g)$ with $R = \norm{f-P_Sf}^2/\norm{P_Sf}^2$ the initial ratio, then
\[
  \norm{H_t(f - P_Sf)}^2 \;\le\; \eta\,\norm{H_tP_Sf}^2
  \quad\text{(\lean{scale\_separation\_time})}.
\]
\end{theorem}

\begin{proof}[Proof sketch]
$\norm{H_tx}^2 = \sum_j e^{-2\lambda_jt}\ip{x}{\varphi_j}^2$ (DYN-4's \lean{dotProduct\_cfc\_mulVec}
with the semigroup law), and $P_Sf$ lives on the modes in $S$ (\lean{modeProj\_dotProduct\_mode}),
where $\lambda_j \le \nu$, which gives the lower bound (\lean{heat\_energy\_ge}). Multiplying the two
bounds, with $e^{-2\mu t} = e^{-2(\mu-\nu)t}e^{-2\nu t}$ and both energies nonnegative, gives the
product form. For the time, $\log(R/\eta) \le 2gt$ gives $e^{-2gt}R \le \eta$
(Mathlib's \lean{Real.log\_le\_iff\_le\_exp}), so $e^{-2gt}\norm{f-P_Sf}^2 \le \eta\norm{P_Sf}^2$, and
the product form divided by $\norm{P_Sf}^2 > 0$ is the claim; when $\norm{f-P_Sf} = 0$ the claim is
trivial.
\end{proof}

With $S$ the first $k$ modes, $\nu = \lambda_k$, $\mu = \lambda_{k+1}$ and $g$ is the eigengap at
the cut. In amplitude rather than energy the rate is $e^{-gt}$, the form in which the question is
usually put. The discrete twin holds under any symmetric operator that multiplies the dropped modes
by at most $\rho_1$ and the kept modes by at least $\rho_2$ in absolute value: after $n$ steps,
$\rho_2^{2n}\norm{N^n(f-P_Sf)}^2\norm{P_Sf}^2 \le \rho_1^{2n}\norm{f-P_Sf}^2\norm{N^nP_Sf}^2$
(\lean{scale\_separation\_pow}, through \lean{pow\_energy\_ge}); for the lazy walk
$\rho_1 = 1-\varepsilon\mu$ and $\rho_2 = 1-\varepsilon\nu$.

\begin{proposition}[the control: no gap, no time]\label{prop:nogap}
On the identity $M = 1$ every eigenvalue is $1$ (\lean{eigenvalues\_one}), so two equal eigenvalues
straddle every cut, and for every $S$, $f$ and $t$,
\[
  \norm{H_t(f-P_Sf)}^2\norm{P_Sf}^2 = \norm{f-P_Sf}^2\norm{H_tP_Sf}^2
  \quad\text{(\lean{no\_gap\_no\_separation})}:
\]
the dropped share of the energy is the same at every time, so if some $t$ makes the dropped part at
most $\eta$ of the kept part, it already was at $t = 0$ (\lean{no\_gap\_no\_time}).
\end{proposition}

\begin{proof}[Proof sketch]
$M\varphi_j = \lambda_j\varphi_j$ with $M = 1$ and $\norm{\varphi_j} = 1$ gives $\lambda_j = 1$. Then
every field's energy under $H_t$ is exactly $e^{-2t}$ times its own (\lean{heat\_energy\_eq}, the
case of a single eigenvalue on every mode a field lives on), and the common factor cancels.
\end{proof}

\subsection{SCL-3: attribution}

\begin{definition}[scores, orderings, the Shapley value]
Let $\alpha$ be a finite set of $n$ sources. A score is $v : 2^\alpha \to K$ over a field $K$ of
characteristic zero: the evaluation of each configuration. An ordering is a bijection
$\sigma : \alpha \simeq \{0, \dots, n-1\}$ (\lean{Palace.Scale.Ordering}); the sources before $i$
are $\operatorname{pred}(\sigma, i) = \{j : \sigma j < \sigma i\}$; the marginal contribution of $i$
to $S$ is $v(S \cup \{i\}) - v(S)$ (\lean{marginal}). The Shapley value is the average over
every ordering of $i$'s marginal contribution to the sources before it,
\[
  \Sh_i(v) = \frac{1}{n!}\sum_{\sigma}\bigl(v(\operatorname{pred}(\sigma,i) \cup \{i\}) -
  v(\operatorname{pred}(\sigma,i))\bigr)
  \quad\text{(\lean{shapley})}.
\]
An attribution $\psi$ assigns a value per source to every score. It is \emph{efficient} when
$\sum_i\psi_i(v) = v(\alpha) - v(\emptyset)$; \emph{symmetric} when $\psi_i(v) = \psi_j(v)$ whenever
$v(S \cup \{i\}) = v(S \cup \{j\})$ for every $S$ holding neither; \emph{null-respecting} when
$\psi_i(v) = 0$ whenever $v(S \cup \{i\}) = v(S)$ for every $S$; \emph{additive} when
$\psi(v+w) = \psi(v) + \psi(w)$ (\lean{Efficient}, \lean{Symmetric}, \lean{NullPlayer},
\lean{Additive}). The unanimity game of $T$ is $u_T(S) = [T \subseteq S]$, the Harsanyi dividend of
$v$ at $T$ is $c_T(v) = \sum_{S \subseteq T}(-1)^{\abs{T}-\abs{S}}v(S)$, and leave-one-out is
$\loo_i(v) = v(\alpha) - v(\alpha \setminus \{i\})$ (\lean{unanimity}, \lean{harsanyi}, \lean{loo}).
\end{definition}

\begin{theorem}[SCL-3, attribution]\label{thm:scl3}
\sloppy
\begin{enumerate}
\item The Shapley value is efficient (\lean{shapley\_efficient}), symmetric
  (\lean{shapley\_symmetric}), null-respecting (\lean{shapley\_nullPlayer}) and additive
  (\lean{shapley\_additive}).
\item Every score is the combination of the unanimity games with its dividends as coefficients,
  $\sum_T c_T(v)\,u_T(S) = v(S)$ (\lean{sum\_harsanyi\_unanimity}), and any attribution with the
  four axioms is the dividend formula
  \[
    \psi_i(v) = \sum_{T \ni i}\frac{c_T(v)}{\abs{T}}
    \quad\text{(\lean{attribution\_eq\_sum\_harsanyi})},
  \]
  so it is the Shapley value (\lean{eq\_shapley\_of\_axioms}, \lean{shapley\_eq\_sum\_harsanyi}).
\item If $v$ has constant marginals, $v(S \cup \{i\}) - v(S) = a_i$ for every $S \not\ni i$, then
  $\loo_i(v) = \Sh_i(v) = a_i$ (\lean{loo\_eq\_shapley\_of\_marginal\_const}).
\item (The witness.) On two sources with $v(S) = [S \neq \emptyset]$, either source sufficient,
  $\Sh_0(v) = \Sh_1(v) = \tfrac12$ while $\loo_0(v) = \loo_1(v) = 0$, and the leave-one-out values
  do not sum to $v(\alpha) - v(\emptyset)$ (\lean{Witness.loo\_confirms\_nothing}).
\end{enumerate}
\end{theorem}

\begin{proof}[Proof sketch]
(1) Along one ordering the marginals telescope through the chain
$\emptyset = C_0 \subset C_1 \subset \dots \subset C_n = \alpha$ of the sources before each position
(\lean{chain}, \lean{sum\_marginal\_pred}), so they sum to $v(\alpha) - v(\emptyset)$ for every
$\sigma$, and dividing by the number of orderings, $n!$ (Mathlib's \lean{Fintype.card\_equiv}),
gives efficiency. Composing with the swap of $i$ and $j$ is a bijection on orderings
(\lean{swapOrd}); under it the sources before $i$ are the sources before $j$ with $i$ replaced by
$j$ when $i$ was among them (\lean{pred\_swap}), and the symmetry hypothesis makes the two marginals
equal in both cases (\lean{marginal\_swap}), so the two sums agree. A null source's marginals vanish
termwise; additivity is the linearity of the sum.

(2) $\sum_{T \subseteq S}c_T(v) = \sum_{R \subseteq S}v(R)\sum_{R \subseteq T \subseteq S}(-1)^{\abs{T}-\abs{R}}$
after exchanging the sums, and the inner sum is $[S = R]$: it is the alternating sum over the
subsets of $S \setminus R$ (\lean{sum\_powerset\_filter\_neg\_one\_pow}, through Mathlib's
\lean{sum\_powerset\_neg\_one\_pow\_card}). On a scaled unanimity game $c\,u_T$, a source outside $T$
is null, two sources inside $T$ are symmetric, and efficiency gives $c$ in total, so an attribution
with the three axioms is $c/\abs{T}$ on $T$ and $0$ elsewhere (\lean{attribution\_unanimity});
additivity over the decomposition (\lean{additive\_sum}, by induction on the finite sum) gives the
dividend formula, which the Shapley value, having the four axioms, satisfies too.

(3) Every marginal of $i$ in the ordering average is $a_i$, and $\loo_i(v)$ is the marginal to
$\alpha \setminus \{i\}$. (4) By symmetry the two Shapley values are equal and by efficiency they sum
to $1$; each leave-one-out is $1 - 1 = 0$.
\end{proof}

The witness is the shape of the Rayleigh family's \lean{loo\_confirms\_nothing}: a rise in
conductance carried by two redundant new paths is confirmed by neither when each is reverted alone.
Leave-one-out measures \emph{necessity}; the Shapley value measures the \emph{share carried}; they
agree exactly when sources do not interact, and where they disagree the store holds the same evidence
twice.

\subsection{SCL-4: power}\label{sec:scl4}

The number of independent resolved units an effect needs is not formalized; it is the normal
approximation, in the form the ledger's \lean{required\_clusters} computes. Let $d_k$ be the mean
paired difference between two arms over the pairs of cluster $k$, for $K$ independent clusters,
with mean $\Delta$ and standard deviation $s$ across clusters. The statistic is the cluster mean
$\bar d \approx \mathcal N(\Delta, s^2/K)$. The two-sided test at level $\alpha$ rejects $\Delta = 0$
when $\abs{\bar d} > z_{1-\alpha/2}\,s/\sqrt K$. At a true effect $\Delta > 0$, up to the negligible
other tail, the rejection probability is $P(Z > z_{1-\alpha/2} - \Delta\sqrt K/s)$, which is at least
$1 - \beta$ exactly when $\Delta\sqrt K/s \ge z_{1-\alpha/2} + z_{1-\beta}$, that is
\begin{equation}\label{eq:power}
  K \;\ge\; \Bigl(\frac{(z_{1-\alpha/2} + z_{1-\beta})\,s}{\Delta}\Bigr)^2,
\end{equation}
rounded up, with $s$ read blinded from a pilot that reports no mean. At $s = 0.08$, $\Delta = 0.02$,
$\alpha = 0.05$ and power $0.8$, $K = 126$; at $\Delta = 0.01$, $503$.

SCL-1 adds a constraint on $K$. The units must be independent, and questions registered within one
relaxation time of each other, about a store whose state has not relaxed between them, are one
observation: a window of $T_c$ cuts yields at most about $T_c\gamma/2$ independent units of a macro
observable (\Cref{prop:sampled}). The cluster rule, one cluster per identity or pair, catches the
dependence within a cut; the relaxation time is the rule for dependence across cuts, so the loop's
rounds are spaced at least a relaxation time apart.

\section{Concepts and intuition}

\paragraph{Why macro needs a gap and a window.}
A macro variable is a coefficient on a slow mode, and it is legitimate as a description exactly
when the fast modes have decayed relative to it. SCL-2 makes the two words precise: ``fast'' and
``slow'' are the eigenvalues on either side of the cut, and the time at which the description holds
is $\log(R/\eta)/(2g)$, so a spectrum with no gap near the cut gives no time at all
(\Cref{prop:nogap}). Measuring the macro variable is a different matter: the estimate from a
window is a time average, and SCL-1 says that it converges only at the rate the gap sets, so the
window must be many relaxation times long whatever its sampling rate. A reader who samples a slowly
mixing store a thousand times in one relaxation time has one reading, not a thousand. The two
thresholds are about different gaps: SCL-1's $\gamma$ is the gap above zero, the rate at which any
mean-zero disturbance averages out; SCL-2's $g$ is the eigengap at the cut, the rate at which the
modes the description drops separate from the ones it keeps.

\paragraph{Why the spectrum must have concentrated first.}
Both thresholds take $\gamma$ and $g$ as given, but on a small sample of a growing corpus they are
noise. The concentration results say when they stop being noise: once the minimum expected degree
exceeds a constant times $\log N$, the eigenvalues of the normalized Laplacian sit within
$O(\sqrt{\log N/d_{\min}})$ of their expectations, so a gap is a property of the corpus once that
fluctuation scale is below it, $d_{\min} \gg \log N/\gamma^2$. The order is therefore: concentrate,
then separate, then average.

\paragraph{Why attribution needs a common currency and power.}
Two configurations can only be compared on the same questions in the same rounds, and the
comparison must be in a unit that adds: the log-score advantage of an arm over the reference is
additive over questions, its expectation under a truth $q$ is $\KL(q\|r) - \KL(q\|p)$, positive
exactly when the arm's forecasts are closer to the truth than the base rate, and its cumulative form
is the log growth of a proportional stake against the reference, so the attribution, the growth
readout and the decision rule read one number. Attribution then needs a value function on
configurations, which is what the currency provides, and a bar: \eqref{eq:power} says how many
independent resolved units a difference of $\Delta$ needs before it can be read, and SCL-1 says how
many such units a window of the store's history can hold.

\paragraph{Why two readings.}
An ablation arm measures leave-one-out, and SCL-3 says that is contribution only when marginals are
constant. The Shapley value is the unique attribution that respects all four axioms, and the witness
is the common case of a store: two sources that say the same thing. Reporting both readings turns
their disagreement into a measurement of redundancy, which is itself something a store should know
about its sources.

\section{The loop}\label{sec:loop}

The loop runs on the prediction ledger of the previous design: an append-only chain whose first
record is the protocol, with forward-only predictions, strictly proper scores, cluster-level
inference, blinding until a pre-registered $N$, and the growth readout. What is added is where the
questions come from (the store itself), what the arms are (configurations), and how the outcome is
attributed.

\begin{algorithm}[h]
\caption{The continuous evaluation loop, one round per cut.}
\begin{algorithmic}[1]
\Require a store with cuts $c_1 < c_2 < \dots$ spaced at least a relaxation time $1/\gamma$ apart; a
  protocol fixing the generators, the arms, the reference arm, the bar $(\Delta, \alpha, 1-\beta)$,
  $N$ and the sampled orderings
\Ensure at $N$ clusters, once: $v(S)$ per arm, $\loo_i$ with its sign-flip interval, $\Sh_i$ with
  its standard error, the growth readout
\For{each cut $c$}
  \State \textbf{register} $Q_c \gets$ generators(store as of $c$), seeded: for each question its
    family, its cluster (the identity or pair it turns on), its close time $c' > c$ and its resolver
  \State \textbf{forecast} for each arm $S$ (the full set, every $\alpha \setminus \{i\}$, the
    configurations of the sampled orderings): $p_S(q)$ for $q \in Q_c$, from the store as of $c$
    under $S$'s grant, recorded with $S$'s pinned digest
  \State \textbf{resolve} every $q$ whose close time is at most $c$, by its resolver over the store
    as of $c$; a prediction at or after the resolution is late and excluded
  \State \textbf{attribute} when the first $N$ clusters of the primary comparison have resolved,
    once: $v(S)$ the mean log-score advantage over the reference; $\loo_i = v(\alpha) -
    v(\alpha\setminus\{i\})$ by the cluster sign-flip test, Holm-corrected; $\Sh_i$ by the average
    of the marginals over the sampled orderings, with $s_i/\sqrt m$
\EndFor
\end{algorithmic}
\end{algorithm}

\paragraph{The generators.}
Four, each a function of the snapshot at the cut and nothing later: \emph{link}, will identity $a$
cite identity $b$ by $c'$, for pairs not yet linked, sampled with the same rule for every arm;
\emph{supersession}, will note $a$ be superseded by $c'$; \emph{version}, will identity $a$ have a
new version by $c'$; \emph{confirmation}, will a source arriving by $c'$ confirm the answer the store
gives at $c$, with the answer's truth condition stated at $c$. Each question's cluster is the
identity or pair it turns on, so two questions about one note are one observation.

\paragraph{The arms.}
A configuration is a subset of the sources, context windows and analysis tools the system may use:
a source is an origin the store reads from, a context window is the depth of history an arm may read
as of, a tool is a verb of the query algebra the arm may call. The arms differ in what they may
read and call, not in who answers; the model and the prompt are in every arm's digest, and a change
to them is a new variant. The reference arm is the family's base rate over the questions resolved
before the cut, which no source informs; the null arm is its copy and must show exactly zero skill.

\paragraph{The Shapley cost.}
With $n$ sources the exact value needs $2^n$ configurations, each an arm forecasting every
question; beyond $n \approx 5$ the permutation estimator samples $m$ orderings at genesis, registers
the $m n$ configurations they run through, and reports $\Sh_i$ with its standard error $s_i/\sqrt m$,
$s_i$ the spread of $i$'s marginal across orderings \cite{castro2009polynomial,ghorbani2019data};
$m$ is fixed before the first prediction, like $N$.

\section{Evaluation}

\paragraph{What is proved.}
The theorems of \Cref{thm:scl1,thm:scl2,thm:scl3} and \Cref{prop:nogap} are machine-checked in
\lean{formal/Palace/Scale.lean} over the model \lean{formal/Palace/Scale/Defs.lean}, on Lean
v4.36.0-rc2 and the Mathlib it pins \cite{moura2021lean4,mathlib2020}, with the sorry budget at zero
and the axioms audited (\lean{propext}, \lean{Classical.choice}, \lean{Quot.sound} only). The
library's export, 55{,}408 declarations from \lean{Init} up, is replayed by Lean's kernel and by
\lean{leanchecker-paranoid}, \lean{lean4lean}, \lean{nanoda}, \lean{con-leche} and \lean{con-ron},
and the 285 statements of the ledger's specification are judged identical to the library's. The
gate is \lean{formal/check.sh}.

\paragraph{The attacks.}
Each family is attacked on an exact shadow over $\mathbb{Q}$ with a negative control
(\lean{formal/PalaceTest/Scale.lean}, seeds 350 to 367, 500 instances each, Plausible). The spectral
gap is an eigenvalue and does not run, so SCL-1 is attacked where the spectrum is rational: complete
bipartite graphs on four nodes with a uniform weight $w$, whose Laplacian has eigenvalues $0$, $wa$,
$wb$ and $w(a+b)$, so the gap is $w\min(a,b)$ exactly and the lazy step $1/(4w)$ keeps every
multiplier in $[0,1]$; the time average of a centered field over windows of up to seven steps is
matrix arithmetic. The window bound and the geometric-sum form survived 495 and 493 instances (the
rest discarded as one-sided bipartitions). SCL-2's discrete twin is attacked on diagonal operators,
whose modes are the coordinates, with the kept set chosen by bits independently of the eigenvalues
so that the cut need not fall in a gap: 500. SCL-3 is attacked on games over three sources with
values from a list, the six orderings enumerated through Mathlib's \lean{Fintype} instance on
$\mathrm{Fin}\,3 \simeq \mathrm{Fin}\,3$: efficiency, additivity, the null source, symmetry, the
Harsanyi decomposition at every configuration, the dividend formula and leave-one-out on an additive
score, 500 each.

The controls, each refuted: the largest eigenvalue in place of the gap, within 5 instances; a field
with a mean, which the walk conserves, within 25; the dropped share ``always falls'', within 19, on a
cut with a kept eigenvalue above a dropped one, where the share rises; the fastest dropped mode's rate
in place of the slowest's, within 24; the equal split, efficient and symmetric, pays a null source,
within 11; leave-one-out is not efficient, within 10 (a score of $1$ on the empty configuration
alone); the uniform average over configurations (the Banzhaf weights), symmetric, null-respecting
and additive, is not efficient, within 8; the dividends without their $1/\abs{T}$, within 6.

\paragraph{What the sibling study measures.}
A separate pull request (\#76) measures the concentration threshold on real data. Its
fetcher retrieves the works of one OpenAlex concept \cite{priem2022openalex} with their dates and
reference lists, resumable by date slice; its measurement, when it lands, is the eigenvalues of the
normalized Laplacian of the growing dated citation graph at successive cuts, and where the gaps stop
moving as a function of the minimum degree. No measurement from it is reported here.

\paragraph{How far below the thresholds a small private corpus is.}
The live store, measured read-only on 2026-11-04 and 2026-11-07 for the previous article: a
backbone of 38 notes and 38 edges, a relation graph of 20 edges on 12 notes, 29 identities with one
version each. With $N = 38$, $\log N \approx 3.6$, mean degree $2$ and minimum degree $1$, the
minimum degree is below $C\log N$ for any $C \ge 1/3$ and the fluctuation scale
$\sqrt{\log N/d_{\min}} \approx 1.9$ exceeds the whole range $[0,2]$ of the normalized Laplacian's
spectrum: every measured gap is noise of the sample. There is one cut, so the window is zero
relaxation times for any $\gamma$; there are no resolved questions, so $K = 0$ against a required
$126$; and one generation, so the supersession and version generators have nothing to ask. The
small corpus is below every threshold by two to four orders of magnitude, and nothing measured on it
about its macro regime is evidence.

\section{Datasets and the recipe}\label{sec:datasets}

The thresholds need an evolving graph with dates, large enough that its spectrum has concentrated
and with a history of many relaxation times; the loop needs events after a cut that a snapshot before
it does not contain. \Cref{tab:datasets} lists five open corpora with their licences, their scale
and what each exercises.

\begin{table}[h]
\centering
\small
\begin{tabular}{p{0.19\linewidth}p{0.16\linewidth}p{0.21\linewidth}p{0.32\linewidth}}
\toprule
corpus & licence & scale & exercises \\
\midrule
OpenAlex works, with \texttt{publication\_date} and \texttt{referenced\_works}
  & CC0 1.0
  & hundreds of millions of works, on the order of a billion dated citation links (the study records what it fetched)
  & a dated citation graph: the link generator at every cut; the gap and the eigengaps of a growing graph against the minimum degree; concept clusters as a partition \\
Wikipedia dumps, articles with revision histories and links
  & CC BY-SA 4.0
  & millions of articles, more than a billion revisions in the English edition
  & version chains: the supersession and version generators; links that appear and are severed; a velocity with many generations \\
commit and issue histories of public repositories
  & each repository's own licence; the platform's terms for the metadata (mirror the metadata only)
  & millions of repositories; hundreds of thousands of commits and issues on a large one
  & commits as supersessions on a path's chain; issues citing commits as a dated relation; many sources for the confirmation generator \\
arXiv metadata through OAI-PMH
  & CC0 1.0 (metadata)
  & more than two million papers, each with dated versions
  & version chains with dates and no citations; a feed the loop registers from, round by round \\
ALFRED, archival economic series with vintages
  & the St.\ Louis Fed's terms of use (most series free with attribution; some third-party series restricted)
  & hundreds of thousands of series, the vintages of the archived ones, decades of them
  & a value behind a fixed identity re-recorded at each vintage; the velocity of a value with its revisions; a confirmation generator whose later source is the next vintage \\
\bottomrule
\end{tabular}
\caption{Open corpora that give an evolving graph with dates, and what each exercises.}
\label{tab:datasets}
\end{table}

Every edge, version or value in them carries a date the corpus assigned, so a cut is a date and a
snapshot as of a cut is reproducible, and the forward-only rule holds against a corpus that already
contains the future only if the arms read snapshots as of the cut and nothing later, which the
instance that holds the corpus enforces by what it lands.

\subsection{The recipe}\label{sec:recipe}

\begin{enumerate}
\item \textbf{Measure.} At a sequence of cuts, the Laplacian's gap $\gamma$ and the eigengaps $g_k$
  near the candidate cuts, for the normalized Laplacian too, with the minimum degree $d_{\min}$ and
  the size $N$, each as a function of the cut.
\item \textbf{Size.} From $\gamma$, the window $T_c \ge 2\sigma^2/(\gamma\delta^2)$ a macro
  observable needs and the effective samples a history of $T_c$ carries; from $g_k$, the time scale
  $1/g_k$ at which a $k$-mode description is legitimate; from $d_{\min}$ and $N$, whether the spectrum
  has concentrated; from the pilot's spread, $K$ by \eqref{eq:power}.
\item \textbf{Grow.} Fetch until the thresholds are exceeded: $d_{\min} \gg \log N/\gamma^2$, a history
  of many relaxation times, enough clusters for $K$. Each fetch is a dated extension, so growing is
  adding cuts, not replacing the corpus.
\item \textbf{Run.} The loop under a protocol written first: the generators, the arms (the full
  system, the ablations, the sampled configurations with $m$ fixed), the primary comparison and the
  bar, the round spacing at a relaxation time.
\item \textbf{Attribute.} At $N$ clusters, once: the sign-flip reading per ablation, the Shapley
  reading with its standard error, the growth readout; where the two readings disagree, name the
  redundancy.
\end{enumerate}

\paragraph{Where it lives.}
The machinery is generic: the generators and resolvers are functions of a store as of a cut,
written over the query algebra and so over any store; the ablation arms are grants; the attribution
readings are functions of the ledger. They extend the ledger's scorer and nothing in the sealed core
imports them. A corpus belongs in an instance of its own, with its own ledger and its own anchored
head; instance identity, an id in the schema and the namespaces derived from it, is what makes a
public corpus of millions of nodes safe to run beside a private store, and the forward-only rule
holds because the instance lands snapshots as of a cut and nothing after.

\section{Limitations}

The proved constant of SCL-1 is the deterministic one, $1/(\gamma T_c)^2$, for the exact diffusion;
the sampled form with the constant $2/(\gamma T)$ and the effective sample size $T\gamma/2$ is derived
for a chain started at stationarity and not machine-checked, and a chain started elsewhere adds a
burn-in of order $\log(N)/\gamma$ that the derivation omits. SCL-2 is stated for a fixed symmetric
operator; a store's graph changes between cuts, and the modes of the graph at one cut are not those
at the next, so the time $1/g$ is a statement about the graph at a cut, read on the next only if the
spectrum has concentrated. The concentration threshold is cited for graphs with independent edges,
which a citation graph is not; its constant and its exponent on real data are the sibling study's to
report, and until it reports nothing here predicts where a measured spectrum stops moving. SCL-3's
witness shows the disagreement between the two readings but proves no bound on it; the general
relation between leave-one-out and the Shapley value depends on the score's interaction structure
and is not characterized here. SCL-4 is the normal approximation at large $K$ and assumes
independent clusters; dependence across clusters, which the sign-flip test's note already records,
makes the effective $K$ smaller than the count. The loop is designed and not built; its generators,
its arms as grants and its attribution readings are proposals to the ledger's design, each with its
falsifier, and none has run. The live store's figures place it below every threshold, which is a
measurement, not a projection of where a larger private store would be.

\section{Future work}

Build the loop over the existing ledger: the four generators and their resolvers as functions of a
store as of a cut, the configuration arms as grants, and the two attribution readings, each
property-tested against the Lean statements before it lands, with the laws entering the ledger as
rows when the code relies on them. Run the sibling study's measurement on the first open corpus and
replace the cited concentration threshold with a measured curve of the eigengaps against the
minimum degree. Formalize the sampled form of SCL-1 over Mathlib's Markov chain library when it
reaches stationary variance bounds, and the burn-in it omits. Characterize when the two attribution
readings agree beyond constant marginals, in terms of the dividends of the score's interaction
terms. Extend the thresholds from a fixed graph to the sequence of graphs, which is the mixing
hypothesis of the previous article and needs a series the store does not have yet.

\section{Conclusion}

A store can study its macro behaviour once its graph has concentrated, once a window of its history
spans many relaxation times, and once the eigengap at the cut it reads separates the modes it keeps
from the ones it drops; each of these is a theorem or a cited theorem with a number attached, and a
small private corpus is below every one of them by orders of magnitude. A store can attribute its
progress to its sources once every configuration forecasts the same questions about its own future,
forward only and pre-registered, in the one currency that adds and that is the growth of a stake
against the base rate; the attribution is then two readings, the ablation's leave-one-out and the
Shapley value, and the proof says they disagree exactly where the store holds the same evidence
twice. The theorems are machine-checked, their attacks survived and their controls refuted; the loop
and the corpora are designed; the first measurement on real data is a sibling study's to report.

\bibliographystyle{plainnat}
\bibliography{refs}
\end{document}
